% from ega4-17-fr.tex
\oldpage[IV-4]{94}

$\inf(Y,Z)$ de $X$ (I, 4.4.3), et l'on a $t=y=z=x$. Au lieu de dire que le couple $(f,g)$ est transversal au point $x$, on dit alors que $Y$ et $Z$ \emph{se coupent transversalement au point $x$} (relativement à $S$). Notant $X_s$, $Y_s$, $Z_s$, $T_s$ les fibres de $X$, $Y$, $Z$, $T$ au point $s$, image de $x$ dans $S$, et tenant compte de (5.2.3), (5.1.9) et (0, 16.5.12), on voit que pour qu'il en soit ainsi (lorsque $X$, $Y$, $Z$ sont lisses sur $S$ au point $x$), il faut et il suffit que $T$ soit lisse sur $S$ au point $x$, et que l'on ait la relation
\[
\label{IV.17.13.7.1}
\tag{17.13.7.1}
\operatorname{codim}_x(T_s,X_s)=\operatorname{codim}_x(Y_s,X_s)+\operatorname{codim}_x(Z_s,X_s).
\]
\end{env}

\begin{proposition}[17.13.8]
\label{IV.17.13.8}
Soient $S$ un préschéma, $X$ un $S$-préschéma localement de présentation finie sur $S$, $Y$, $Z$ deux sous-préschémas de $X$, $T=\inf(Y,Z)$ le sous-préschéma « intersection » de $Y$ et $Z$, $x$ un point de $T$. Les conditions suivantes sont équivalentes:
\begin{enumerate}
  \item[a)] L'injection canonique $i:Y\to X$ est un morphisme transversal à $Z$ au point $x$, relativement à $S$ (17.13.3).
  \item[a')] L'injection canonique $j:Z\to X$ est un morphisme transversal à $Y$ au point $x$, relativement à $S$.
  \item[a'')] $Y$ et $Z$ se coupent transversalement au point $x$, relativement à $S$.
  \item[b)] $X$, $Y$, $Z$ sont lisses sur $S$ au point $x$, et l'homomorphisme $\rho$ (17.13.5.4) est tel que
  \[
  \rho\otimes1:\Omega^1_{X/S}\otimes_{\sh{O}_X}k(x)\to
  (\Omega^1_{Y/S}\otimes_{\sh{O}_Y}k(x))\oplus
  (\Omega^1_{Z/S}\otimes_{\sh{O}_Z}k(x))
  \]
  est injectif.
\end{enumerate}
Lorsque $x$ est rationnel sur $k(s)$ ($s$ image de $x$ dans $S$) ces conditions sont aussi équivalentes à la suivante:
\begin{enumerate}
  \item[b')] L'homomorphisme
  \[
  \mathrm{T}_x(i)-\mathrm{T}_x(j):\mathrm{T}_{Y/S}(x)\oplus\mathrm{T}_{Z/S}(x)\to\mathrm{T}_{X/S}(x)
  \]
  est surjectif.
\end{enumerate}
En outre, lorsque les conditions équivalentes a) à b) sont vérifiées au point $x$, elles le sont dans un voisinage de $x$ dans $T$, et en restreignant $X$ à un voisinage de $x$, la suite (17.13.6.2) est exacte, et l'on a un isomorphisme canonique
\[
\label{IV.17.13.8.1}
\tag{17.13.8.1}
\sh{N}_{T/X}\xrightarrow{\sim}
(\sh{N}_{Y/X}\otimes_{\sh{O}_Y}\sh{O}_T)\oplus
(\sh{N}_{Z/X}\otimes_{\sh{O}_Z}\sh{O}_T).
\]
\end{proposition}

Les conditions a), a'), a'') impliquent toutes que $X$, $Y$, $Z$ sont lisses sur $S$ au point $x$. Soient en outre, $m$, $m-a$, $m-b$ les dimensions relatives de $X$, $Y$, $Z$ sur $S$ au point $x$. Il résulte alors de (17.13.2) appliqué en remplaçant $X'$ par $Z$ et $Y'$ par $T=Y\times_XZ$, que les conditions a), a') équivalent toutes deux à dire que $T$ est lisse sur $S$ au point $x$ et de dimension relative $m-a-b$ en ce point; mais en vertu de (17.13.7), cela signifie précisément que la condition a'') est vérifiée, d'où l'équivalence de a), a') et a''). L'équivalence de a'') et de b) (ou b') lorsque $x$ est rationnel sur $k(s)$) a été prouvée dans (17.13.6), ainsi que le fait que si ces conditions sont satisfaites au point $x$, elles le sont dans un voisinage de $x$, et l'exactitude de la suite (17.13.6.2) dans un tel voisinage. Reste à définir l'isomorphisme canonique (17.13.8.1).

